%!TeX TS-program = lualatex
%!TeX encoding = UTF-8 Unicode
%!TeX spellcheck = en-UK
% -*- coding: UTF-8; -*-
% vim: set fenc=utf-8
%%%-----------------------------------------------------------------
% Standalone figure: chaining two Spiking Motifs in a recurrent HD-SNN.
% Build:  lualatex -interaction=nonstopmode fig_izhikevich.tex
%
% Everything is data-driven: to change the network, edit only the block
% "Motif definition" below.
%   * \Nneurons                      size of the population;
%   * \greenMotif / \brownMotif      comma separated "neuron/delay" pairs;
%   * \greenSyncTime / \brownSyncTime  instant of detection.
% Spike times, trajectories and delay labels are all *derived* from those
% lists, so the wiring diagram (left) and the raster (right) cannot drift
% apart: the delay printed on an arrow is exactly the horizontal distance
% covered by its spike, and an input spike is drawn where delay + time =
% instant of detection.
%%%-----------------------------------------------------------------
\documentclass[tikz,border=2pt]{standalone}     % "border" adds optional margin
% -------------------------------------------------------------
\usepackage{amsmath,amssymb}            % math, symbols
\usepackage{tikz}
\usetikzlibrary{arrows,calc,positioning,shapes,decorations.pathreplacing}
\usepackage{pgfplots}                   % if you used pgfplots
\pgfplotsset{compat=newest}             % optional - keep the same version as before
% -------------------------------------------------------------

%==================== Motif definition ==========================
% --- Population ---------------------------------------------------------
\def\Nneurons{8}         % neurons in the recurrent population (1 ... \Nneurons)
\def\Nmotif{8}           % highest index used by the two motifs; above: background

% --- Motifs: "presynaptic neuron / axonal delay in ms" ------------------
\def\greenMotif{8/1, 3/6, 4/2, 7/8}     % a3, a4, a7, a8 -> b5
\def\brownMotif{8/3, 3/8, 4/4, 5/2}     % a3, a4, a5, a8 -> b2

% --- Neuron emitting the detected (polychronous) spike ------------------
\def\greenTarget{5}
\def\brownTarget{2}

% --- Instant of detection (ms) ------------------------------------------
% The green spike re-enters the population and feeds the dark-red motif,
% which is what chains the two motifs.
\def\greenSyncTime{9}
\def\brownSyncTime{11}

% --- Background activity of the unconnected neurons, "time/neuron" ------
\def\spontaneousSpikes{1/1, 1/4, 3/4, 5/7, 6/2, 12/2}

% Where the delay labels sit along their arrow, in [0, 1]. The two motifs are
% pulled apart so that their labels do not overlap where the bundles cross.
\def\greenLabelPos{0.62}
\def\brownLabelPos{0.38}

% --- Geometry -----------------------------------------------------------
\def\preLabelX{-0.15}    % abscissa of the "pre" title
\def\postLabelX{4.15}    % abscissa of the "post" title
\def\bRightShift{4.55}   % abscissa of the post-synaptic column
\def\panelShift{7.4cm}   % gap between the wiring diagram and the raster
\def\spikeTop{0.75}      % height of a spike, in row units

% Derived quantities. Every one is evaluated by pgfmath, never left as text.
\pgfmathtruncatemacro\lastRow{\Nneurons - 1}
\pgfmathtruncatemacro\lastTick{\brownSyncTime + 1}
\pgfmathsetmacro\rowGreen{\Nneurons - \greenTarget}
\pgfmathsetmacro\rowBrown{\Nneurons - \brownTarget}
\pgfmathsetmacro\rasterRight{\lastTick + 0.8}
\pgfmathsetmacro\axisRight{\lastTick + 0.4}
\pgfmathsetmacro\windowTop{\lastRow + \spikeTop + 0.3}
\pgfmathsetmacro\titleY{\Nneurons - 0.45}

%======================= Styles and helpers ======================
\tikzset{
    > = stealth,
    neuron/.style={circle, draw=blue!70!black, line width=1.1pt,
        fill=blue!8, minimum size=0.95cm, inner sep=0pt,
        font=\sffamily\small},
    % Neighbours taking no part in the two motifs are faded out.
    neuron background/.style={neuron, draw=blue!70!black!35, fill=blue!3,
        text=blue!70!black!45},
    panelbox/.style={rounded corners=2pt, line width=1.1pt,
        dash pattern=on 2.2pt off 3.2pt},
    layer/.style={draw=blue!60!black, line width=1.2pt},
    layer background/.style={draw=blue!60!black!30, line width=1pt},
    brownlink/.style={draw=black!60!red!75!brown, line width=1.7pt},
    greenlink/.style={draw=green!60!black, line width=1.7pt},
    trigger/.style={draw=orange!90!black, line width=3.5pt},
    greenout/.style={draw=green!45!black, line width=3.5pt},
    redout/.style={draw=red!90!black, line width=4.4pt},
    spontaneous/.style={draw=black!45, line width=2.2pt, opacity=0.4},
    recurrent/.style={densely dotted, draw=blue!55!black!35, line width=0.9pt},
    nodeLabel/.style={sloped, above, fill=white, inner sep=2pt,
        font=\sffamily\scriptsize},
    panelTitle/.style={font=\itshape\Large},
    axisLabel/.style={font=\Large},
    tick/.style={font=\small},
}

% Vertical spike mark of height \spikeTop standing on the row of neuron #2.
% #1 = time, #2 = neuron index, #3 = style.
\newcommand\spikeMark[3]{%
   \pgfmathsetmacro\row{\Nneurons - #2}%
   \draw[#3] (#1,\row) -- (#1,{\row + \spikeTop});%
}

% Wiring diagram: one arrow per input spike, converging on the target.
% #1 = style, #2 = motif list, #3 = target neuron, #4 = label position.
\newcommand\drawWiring[4]{%
   \pgfmathsetmacro\targetRow{\Nneurons - #3}%
   \foreach \i/\d in #2 {%
     \pgfmathsetmacro\row{\Nneurons - \i}%
     \draw[#1] (0,\row) -- node[nodeLabel, pos=#4] {$\d$\,ms} (\bRightShift,\targetRow);%
   }%
}

% Raster: one trajectory per input spike, converging on the detected spike.
% #1 = style, #2 = motif list, #3 = instant of detection, #4 = target neuron.
\newcommand\drawTrajectories[4]{%
   \pgfmathsetmacro\targetRow{\Nneurons - #4}%
   \foreach \i/\d in #2 {%
     \pgfmathsetmacro\spikeTime{#3 - \d}%
     \pgfmathsetmacro\row{\Nneurons - \i}%
     \draw[#1] (\spikeTime,\row) -- (#3,\targetRow);%
   }%
}

%------------------- START OF THE PICTURE -------------------
\begin{document}
\begin{tikzpicture}[
    x=1.2cm,
    y=1.2cm,
    line cap=round,
    line join=round,
    font=\sffamily,
]

%=== Left panel: the two motifs, seen from the wiring diagram ============
\node[panelTitle, anchor=south] at (\preLabelX,\titleY) {pre};
\node[panelTitle, anchor=south west] at (\postLabelX,\titleY) {post};

% The population is recurrent: neuron \i has a pre- and a post-synaptic side.
\foreach \i in {1,...,\Nneurons} {%
   \pgfmathsetmacro\row{\Nneurons - \i}%
   \coordinate (a\i) at (0,\row);
   \coordinate (b\i) at (\bRightShift,\row);
}

% Recurrent link between the two sides of the same neuron. Drawn with a neutral
% style for every neuron: motif membership is carried by the colored bundles,
% not by this link.
\foreach \i in {1,...,\Nneurons} {%
    \draw[recurrent] (a\i) -- (b\i);
}

\drawWiring{brownlink}{\brownMotif}{\brownTarget}{\brownLabelPos}
\drawWiring{greenlink}{\greenMotif}{\greenTarget}{\greenLabelPos}

% Nodes on top of the links; motif neurons at full contrast, others faded.
\foreach \i in {1,...,\Nneurons} {%
   \ifnum\i>\Nmotif\relax
     \node[neuron background] at (a\i) {$a_{\i}$};
     \node[neuron background] at (b\i) {$b_{\i}$};
   \else
     \node[neuron] at (a\i) {$a_{\i}$};
     \node[neuron] at (b\i) {$b_{\i}$};
   \fi
}

%=== Right panel: the same motifs, seen from the spike train =============
\begin{scope}[xshift=\panelShift]

   % Context windows, first, so that spikes stay in the foreground.
   \draw[panelbox, draw=green!45!black, opacity=.35]
     (0.15,-0.12) rectangle (\greenSyncTime,\windowTop);
   \draw[panelbox, draw=black!35!red, opacity=.35]
     (2.0,-0.10) rectangle (\brownSyncTime,\windowTop);

   % Membrane traces.
   \foreach \r in {0,...,\lastRow} {%
     \pgfmathtruncatemacro\neuronAtRow{\Nneurons - \r}%
     \ifnum\neuronAtRow>\Nmotif\relax
       \draw[layer background] (0,\r) -- (\rasterRight,\r);
     \else
       \draw[layer] (0,\r) -- (\rasterRight,\r);
     \fi
   }

   % Axonal trajectories.
   \drawTrajectories{greenlink}{\greenMotif}{\greenSyncTime}{\greenTarget}
   \drawTrajectories{brownlink}{\brownMotif}{\brownSyncTime}{\brownTarget}

   % Trigger spikes feeding the green motif.
   \foreach \i/\d in \greenMotif {%
     \pgfmathsetmacro\spikeTime{\greenSyncTime - \d}%
     \spikeMark{\spikeTime}{\i}{trigger}%
   }

   % Trigger spikes feeding the dark-red motif. The spike that re-enters the
   % network is the green output spike, so it is not redrawn in orange.
   \foreach \i/\d in \brownMotif {%
     \pgfmathsetmacro\spikeTime{\brownSyncTime - \d}%
     \pgfmathtruncatemacro\isReentry{int((\i == \greenTarget) * (\spikeTime == \greenSyncTime))}%
     \ifnum\isReentry=0\relax
       \spikeMark{\spikeTime}{\i}{trigger}%
     \fi
   }

   % Background activity of the unconnected neurons.
   \foreach \t/\i in \spontaneousSpikes {%
     \spikeMark{\t}{\i}{spontaneous}%
   }

   % Detected spikes: the green one, then the dark-red one it triggers.
   \spikeMark{\greenSyncTime}{\greenTarget}{greenout}
   \spikeMark{\brownSyncTime}{\brownTarget}{redout}

   % Time axis.
   \draw[->, line width=1.3pt] (0.0,-0.68) -- (\axisRight,-0.68)
     node[above, axisLabel] {time (ms)};
   \foreach \t in {0,...,\lastTick} {
     \draw (\t,-0.78) -- +(0,-0.05) node[below, tick] {$\t$};
   }

   % Countdown axis: time left before the green motif is detected.
   \foreach \t in {0,...,\greenSyncTime} {
     \pgfmathsetmacro\xpos{\greenSyncTime - \t}%
     \draw[green!60!black] (\xpos,-0.24) -- +(0,-0.05)
       node[below, tick] {$\t$};
   }

\end{scope}
\end{tikzpicture}
%------------------- END OF THE PICTURE -----------------
\end{document}
%================================================================
